\documentclass[border=10pt]{standalone}
\usepackage[no-math]{fontspec}
\usepackage[brazil]{babel}
\usepackage{amsmath,amssymb}
\IfFileExists{C:/Windows/Fonts/Lato_Regular.ttf}{%
  \setmainfont{Lato}[
    Path = C:/Windows/Fonts/,
    UprightFont = Lato_Regular.ttf,
    BoldFont = Lato_Bold.ttf,
    ItalicFont = Lato_Italic.ttf,
    BoldItalicFont = Lato_BoldItalic.ttf
  ]
  \setsansfont{Lato}[
    Path = C:/Windows/Fonts/,
    UprightFont = Lato_Regular.ttf,
    BoldFont = Lato_Bold.ttf,
    ItalicFont = Lato_Italic.ttf,
    BoldItalicFont = Lato_BoldItalic.ttf
  ]
}{%
  \RequirePackage[default]{lato}
}
\renewcommand{\familydefault}{\sfdefault}

\usepackage{xcolor}
\usepackage{tikz}
\usepackage{pgfplots}
\pgfplotsset{compat=1.18}
\usetikzlibrary{arrows.meta,positioning,shapes.geometric,fit,backgrounds,calc,decorations.pathreplacing,patterns}

% Paleta institucional UEFS / Bioengenharia de Solos
\definecolor{bsNavy}{HTML}{1F4268}
\definecolor{bsNavyDark}{HTML}{152E4A}
\definecolor{bsGold}{HTML}{B38A3D}
\definecolor{bsGoldDark}{HTML}{8C6927}
\definecolor{bsTeal}{HTML}{0E7C7B}
\definecolor{bsForest}{HTML}{2D6A4F}
\definecolor{bsAlert}{HTML}{9E2A2B}
\definecolor{bsEarth}{HTML}{8C5A3C}
\definecolor{bsWater}{HTML}{2E86AB}
\definecolor{bsInk}{HTML}{26313D}
\definecolor{bsMuted}{HTML}{5B6573}

\definecolor{bgNavy}{HTML}{EDF2F7}
\definecolor{bgGold}{HTML}{F8F3E9}
\definecolor{bgTeal}{HTML}{EBF6F5}
\definecolor{bgForest}{HTML}{EEF5F0}
\definecolor{bgAlert}{HTML}{FDF2F2}

\tikzset{
  >=Stealth,
  bsBox/.style={
    draw=#1, fill=white, rounded corners=5pt, line width=1.15pt,
    inner sep=7pt, align=center, font=\footnotesize\sffamily, text=bsInk
  },
  bsTint/.style={
    draw=#1, fill=#1!8!white, rounded corners=5pt, line width=1.15pt,
    inner sep=7pt, align=center, font=\footnotesize\sffamily, text=bsInk
  },
  bsBadge/.style={
    fill=#1, text=white, rounded corners=9pt, inner xsep=8pt, inner ysep=3pt,
    font=\scriptsize\bfseries\sffamily
  },
  bsTitle/.style={font=\normalsize\bfseries\sffamily, text=bsNavy, align=center},
  bsSub/.style={font=\footnotesize\sffamily, text=bsMuted, align=center},
  bsArrow/.style={->, >=Stealth, line width=1.35pt},
  bsDArrow/.style={<->, >=Stealth, line width=1.35pt}
}


\begin{document}
\begin{tikzpicture}
  \hyphenpenalty=10000\exhyphenpenalty=10000

  % Título
  \node[bsTitle] at (2.9, 5.75)
    {Canaletas Vegetadas (Bioswales): Tipologia e Mecanismo de Filtração e Infiltração};
  \node[bsSub] at (2.9, 5.28)
    {Seção transversal típica com taludes suaves e contraste hidrológico frente à drenagem convencional};

  % ============================================================
  % PAINEL ESQUERDO: SEÇÃO TRANSVERSAL DO BIOSWALE
  % ============================================================
  \node[bsBox=bsWater!60!white, minimum width=9.5cm, minimum height=7.15cm, anchor=north west] (leftp) at (-4.35, 4.70) {};
  \node[bsBadge=bsWater] at (leftp.north) {SEÇÃO TRANSVERSAL DE CANALETA VEGETADA};

  \begin{scope}[shift={(-3.8, -0.75)}]
    % Solo natural ao fundo
    \fill[bsEarth!16!white]
      (0, 0.15) -- (8.4, 0.15) -- (8.4, 3.10) -- (6.3, 3.10)
      .. controls (5.4, 2.25) and (5.0, 2.25) .. (4.2, 2.25)
      .. controls (3.4, 2.25) and (3.0, 2.25) .. (2.1, 3.10)
      -- (0, 3.10) -- cycle;

    % Substrato poroso drenante (meio filtrante)
    \fill[bsGold!28!white]
      (2.4, 2.60) .. controls (3.2, 2.25) and (5.2, 2.25) .. (6.0, 2.60)
      -- (6.0, 1.38) -- (2.4, 1.38) -- cycle;
    \draw[bsGoldDark, dashed, line width=0.8pt] (2.4, 1.38) -- (6.0, 1.38);

    % Camada de brita + geotêxtil
    \fill[bsMuted!25!white] (2.4, 0.70) rectangle (6.0, 1.38);
    \draw[bsInk, dotted, line width=1pt] (2.4, 1.38) -- (2.4, 0.70) -- (6.0, 0.70) -- (6.0, 1.38);

    % Dreno subsuperficial perfurado
    \fill[white, draw=bsNavy, line width=1.4pt] (4.2, 1.03) circle (0.24);
    \fill[bsWater!35!white] (4.0, 0.95) arc (200:340:0.21) -- cycle;
    \draw[bsNavy, dashed, line width=0.8pt] (4.2, 1.03) circle (0.16);

    % Lâmina d'água temporária
    \fill[bsWater!25!white]
      (2.45, 2.68) -- (5.95, 2.68)
      .. controls (5.2, 2.25) and (3.2, 2.25) .. (2.45, 2.68) -- cycle;
    \draw[bsWater, line width=1.1pt] (2.45, 2.68) -- (5.95, 2.68);

    % Perfil da superfície côncava
    \draw[bsEarth, line width=1.6pt]
      (0.1, 3.10) -- (2.1, 3.10)
      .. controls (3.0, 2.25) and (3.4, 2.25) .. (4.2, 2.25)
      .. controls (5.0, 2.25) and (5.4, 2.25) .. (6.3, 3.10) -- (8.3, 3.10);

    % Gramíneas e raízes (Manning n elevado)
    \foreach \x/\y in {2.5/2.72, 2.9/2.46, 3.3/2.32, 3.7/2.26, 4.1/2.25, 4.5/2.26, 4.9/2.32, 5.3/2.46, 5.7/2.72} {
      \draw[bsForest, line width=1.1pt] (\x, \y) -- (\x-0.06, \y+0.34);
      \draw[bsForest, line width=1.1pt] (\x, \y) -- (\x+0.07, \y+0.38);
      \draw[bsEarth!80!black, line width=0.6pt] (\x, \y) -- (\x-0.05, \y-0.34);
      \draw[bsEarth!80!black, line width=0.6pt] (\x, \y) -- (\x+0.06, \y-0.37);
    }

    % Setas de processos
    \draw[->, >=Stealth, bsWater, line width=1.15pt] (3.25, 2.20) -- (3.25, 1.54);
    \draw[->, >=Stealth, bsWater, line width=1.15pt] (5.15, 2.20) -- (5.15, 1.54);
    \draw[->, >=Stealth, bsTeal, line width=1.1pt] (4.2, 2.05) -- (4.2, 1.36);

    % Rótulos técnicos bem separados
    \node[font=\scriptsize\bfseries, text=bsInk, align=center, fill=white, rounded corners=2pt, inner sep=2pt]
      at (1.15, 3.85) {Talude suave\\(3H:1V--4H:1V)};
    \draw[->, >=Stealth, bsInk, line width=0.8pt] (1.65, 3.55) -- (2.35, 2.85);

    \node[font=\scriptsize\bfseries, text=bsForest, align=center]
      at (4.95, 3.85) {Vegetação ($\uparrow n$ de Manning)\\e sedimentação superficial};
    \draw[->, >=Stealth, bsForest, line width=0.8pt] (4.65, 3.50) -- (4.35, 3.05);

    \node[font=\scriptsize\bfseries, text=bsWater, align=left]
      at (1.15, 1.85) {Infiltração\\e biofiltração};
    \draw[->, >=Stealth, bsWater, line width=0.9pt] (2.00, 1.85) -- (3.12, 1.85);

    \node[font=\scriptsize\bfseries, text=bsGoldDark, align=left]
      at (7.28, 1.92) {Substrato poroso\\(areia + composto)};
    \draw[->, >=Stealth, bsGoldDark, line width=0.9pt] (6.15, 1.92) -- (5.75, 1.92);

    \node[font=\scriptsize\bfseries, text=bsNavy, align=left]
      at (7.28, 1.00) {Dreno perfurado\\+ brita + geotêxtil};
    \draw[->, >=Stealth, bsNavy, line width=0.9pt] (6.15, 1.00) -- (4.55, 1.00);

    \node[font=\scriptsize\itshape, text=bsMuted] at (1.1, 0.48) {Solo natural};
  \end{scope}

  % ============================================================
  % PAINEL DIREITO: CONCRETO VS. CANALETA VERDE (SbN)
  % ============================================================
  \node[bsTint=bsAlert, text width=5.15cm, minimum height=3.25cm, anchor=north west] (comp1) at (5.45, 4.70) {
    \vspace{0.25cm}
    \textbf{\color{bsAlert}Canaleta de Concreto (Cinza)}\\[4pt]
    \scriptsize
    \textbullet\ Baixa rugosidade ($n \approx 0{,}013$) e alta velocidade ($V \uparrow$)\\[2pt]
    \textbullet\ Escoamento 100\% superficial, sem recarga do lençol\\[2pt]
    \textbullet\ Concentra e antecipa picos de cheia a jusante\\[2pt]
    \textbullet\ Transfere sedimentos e poluentes difusos
  };
  \node[bsBadge=bsAlert] at (comp1.north) {DRENAGEM CONVENCIONAL};

  \node[bsTint=bsTeal, text width=5.15cm, minimum height=3.25cm, anchor=north west] (comp2) at (5.45, 0.80) {
    \vspace{0.25cm}
    \textbf{\color{bsTeal}Canaleta Verde / SbN}\\[4pt]
    \scriptsize
    \textbullet\ Alta rugosidade ($n \approx 0{,}04\text{--}0{,}15$) e dissipação de energia\\[2pt]
    \textbullet\ Amortecimento de vazão de pico e aumento do tempo de concentração\\[2pt]
    \textbullet\ Infiltração, retenção de carga difusa e recarga hídrica
  };
  \node[bsBadge=bsTeal] at (comp2.north) {BIOENGENHARIA / SbN};

\end{tikzpicture}
\end{document}
